\documentclass[11pt,a4paper]{article}

% ============================================================
%  QIMMA -- COURS : Dérivation et étude des fonctions
%  2ème Bac Sciences Mathématiques (A et B)
%  Standard exhaustif : définitions formelles + démonstrations
%  Basé sur templates/qimma-template.tex (charte Qimma)
% ============================================================

\usepackage[french]{babel}
\usepackage[utf8]{inputenc}
\usepackage[T1]{fontenc}
\usepackage{lmodern}
\usepackage[a4paper,margin=2cm,top=2.3cm,bottom=2.6cm]{geometry}
\usepackage{amsmath,amssymb}
\usepackage{xcolor}
\usepackage{graphicx}
\usepackage{tikz}
\usetikzlibrary{shadows,calc}
\usepackage[most]{tcolorbox}
\usepackage{titlesec}
\usepackage{fancyhdr}
\usepackage{enumitem}
\usepackage{pifont}
\usepackage{array}
\usepackage{colortbl}
\usepackage{hyperref}

% ------------------- CHARTE QIMMA -------------------
\definecolor{qteal}{HTML}{0d9488}
\definecolor{qtealdark}{HTML}{134e4a}
\definecolor{qamber}{HTML}{fbbf24}
\definecolor{qambertext}{HTML}{92400e}
\definecolor{ink}{HTML}{1f2937}
\definecolor{inkgray}{HTML}{6b7280}
\definecolor{tealpale}{HTML}{e6f5f3}
\colorlet{colDef}{qteal}
\definecolor{colProp}{HTML}{0f766e}
\definecolor{colEx}{HTML}{b45309}
\definecolor{colRem}{HTML}{9d174d}
\color{ink}

\graphicspath{{../../../../assets/}}
\newcommand{\logofile}{../../../../assets/qimma-logo.pdf}

% ------------------- METADONNEES -------------------
\newcommand{\matiere}{Mathématiques}
\newcommand{\niveau}{2\up{ème} Bac -- Sciences Mathématiques}
\newcommand{\chapitretitre}{Dérivation et étude des fonctions}

% ------------------- EN-TETE / PIED DE PAGE -------------------
\pagestyle{fancy}
\fancyhf{}
\renewcommand{\headrulewidth}{0pt}
\renewcommand{\footrulewidth}{0.5pt}
\fancyfoot[L]{\raisebox{-0.15cm}{\includegraphics[height=0.35cm]{\logofile}}~\small\sffamily\color{inkgray}Qimma}
\fancyfoot[C]{\small\sffamily\color{inkgray}\matiere\ -- \niveau}
\fancyfoot[R]{\small\sffamily\color{qtealdark}\thepage}

% ------------------- TITRES DE SECTION -------------------
\titleformat{\section}
  {\normalfont\sffamily\Large\bfseries\color{qtealdark}}
  {\tikz[baseline=-3.2pt]{\node[circle,fill=qteal,text=white,inner sep=0pt,minimum size=8mm]{\sffamily\bfseries\thesection};}}
  {10pt}{}
  [\vspace{-2mm}{\color{qamber}\titlerule[1.6pt]}\vspace{2mm}]
\titlespacing*{\section}{0pt}{16pt}{10pt}

\titleformat{\subsection}
  {\normalfont\sffamily\large\bfseries\color{qtealdark}}
  {\color{qteal}\ding{212}}{6pt}{}
\titlespacing*{\subsection}{0pt}{12pt}{6pt}

% ------------------- BOITES PEDAGOGIQUES -------------------
\newtcolorbox{fiche}[3][]{
  enhanced, breakable,
  colback=white, colframe=#2,
  boxrule=1pt, arc=2mm,
  top=7mm, left=4mm, right=4mm, bottom=3mm,
  drop shadow={black!25, shadow xshift=1mm, shadow yshift=-1mm},
  attach boxed title to top left={xshift=4mm, yshift=-3mm, yshifttext=-1mm},
  boxed title style={colback=#2, arc=1.5mm, boxrule=0pt, left=2mm, right=2mm},
  coltitle=white, fonttitle=\sffamily\bfseries,
  title={#3}, #1
}
\newenvironment{definition}[1][Définition]
  {\begin{fiche}{colDef}{\ding{110}~~#1}}{\end{fiche}}
\newenvironment{propriete}[1][Propriété]
  {\begin{fiche}{colProp}{\ding{72}~~#1}}{\end{fiche}}
\newenvironment{exemple}[1][Exemple]
  {\begin{fiche}{colEx}{\ding{217}~~#1}}{\end{fiche}}
\newenvironment{remarque}[1][Remarque]
  {\begin{fiche}{colRem}{\ding{43}~~#1}}{\end{fiche}}

% Démonstration (hors boîte, style discret)
\newenvironment{demo}
  {\par\smallskip\noindent{\sffamily\bfseries\color{qtealdark}Démonstration.}\ \itshape}
  {\ \normalfont\hfill$\blacksquare$\par\medskip}

\hypersetup{colorlinks=true, linkcolor=qtealdark, urlcolor=qtealdark}

% ============================================================
\begin{document}

% ------------------- BANNIERE DE TITRE -------------------
\begin{tcolorbox}[
  enhanced, boxrule=0pt, arc=4mm,
  interior style={left color=qtealdark, right color=qteal},
  top=8mm, bottom=8mm, left=10mm, right=32mm,
  drop shadow={black!35, shadow xshift=1.5mm, shadow yshift=-1.5mm},
  before skip=0pt, after skip=9mm,
  overlay={
    \node[anchor=north west] at ([xshift=6mm,yshift=-6mm]frame.north west)
      {\includegraphics[height=1.25cm]{\logofile}};
    \node[anchor=north east, fill=qamber, text=qambertext, rounded corners=2pt,
          inner sep=4pt, font=\sffamily\bfseries\small] at ([xshift=-6mm,yshift=-6mm]frame.north east) {COURS};
  }
]
\hspace{1.6cm}\centering
{\sffamily\bfseries\color{white}\fontsize{23}{27}\selectfont \chapitretitre}\\[3mm]
{\sffamily\color{white!90}\large \matiere\ \textbullet\ \niveau}
\end{tcolorbox}

% ------------------- OBJECTIFS -------------------
\begin{tcolorbox}[
  enhanced, colback=tealpale, colframe=qteal, boxrule=0.8pt, arc=2mm,
  left=5mm, right=5mm, top=3mm, bottom=3mm,
  title={\sffamily\bfseries\color{qtealdark}\ding{51}~~Objectifs du chapitre},
  coltitle=qtealdark, colbacktitle=tealpale, boxrule=0.8pt,
  attach title to upper={\par\vspace{1mm}}
]
\begin{itemize}[leftmargin=6mm, label=\ding{212}, itemsep=1mm]
  \item Étudier la dérivabilité d'une fonction en un point (à droite, à gauche) et interpréter géométriquement (tangente, point anguleux, tangente verticale).
  \item Maîtriser toutes les formules de dérivation : opérations, composées usuelles, dérivée d'une fonction réciproque, puissances rationnelles.
  \item Démontrer et appliquer le théorème de Rolle, le théorème des accroissements finis (TAF) et son inégalité.
  \item Relier rigoureusement signe de la dérivée et monotonie, déterminer et justifier les extremums locaux.
  \item Étudier la concavité, les points d'inflexion, classifier toutes les branches infinies, et mener l'étude complète d'une fonction jusqu'au tracé.
  \item Résoudre les exercices type examen national : tangentes, position relative, encadrements par convexité, suites liées à une fonction.
\end{itemize}
\end{tcolorbox}

\vspace{4mm}

% ------------------- SECTION 1 -------------------
\section{Dérivabilité en un point}

\begin{definition}[Nombre dérivé -- définition formelle]
Soit $f$ définie sur un intervalle ouvert contenant $a$. On dit que $f$ est \textbf{dérivable en $a$} lorsque le taux d'accroissement $\tau(x) = \dfrac{f(x) - f(a)}{x - a}$ admet une limite \textbf{finie} quand $x \to a$ :
\[
  \lim_{x \to a} \frac{f(x) - f(a)}{x - a} = \ell \in \mathbb{R}.
\]
Ce nombre $\ell$, noté $f'(a)$, est le \textbf{nombre dérivé} de $f$ en $a$. De façon équivalente, en posant $h = x - a$ :
\[
  f'(a) = \lim_{h \to 0} \frac{f(a + h) - f(a)}{h}.
\]
On définit de même la \textbf{dérivabilité à droite} ($x \to a^+$, notée $f'_d(a)$) et \textbf{à gauche} ($x \to a^-$, notée $f'_g(a)$), en restreignant la limite à $x > a$ ou $x < a$.
\end{definition}

\begin{propriete}
\begin{itemize}[leftmargin=6mm, itemsep=0.5mm]
  \item $f$ est dérivable en $a$ si et seulement si $f'_d(a)$ et $f'_g(a)$ existent et sont \textbf{égales} (auquel cas $f'(a)$ est leur valeur commune).
  \item \textbf{Interprétation géométrique} : si $f$ est dérivable en $a$, la courbe $\mathcal{C}_f$ admet au point $A\bigl(a\,;f(a)\bigr)$ une \textbf{tangente} de coefficient directeur $f'(a)$, d'équation
  \[
    y = f'(a)\,(x - a) + f(a).
  \]
  \item Si $f'_d(a) \neq f'_g(a)$ (toutes deux finies), la courbe admet un \textbf{point anguleux} en $A$ : deux demi-tangentes de pentes différentes.
  \item Si le taux d'accroissement tend vers $+\infty$ ou $-\infty$ (à droite et/ou à gauche), la courbe admet une \textbf{demi-tangente verticale} en $A$, d'équation $x = a$.
  \item \textbf{Dérivable $\Rightarrow$ continue.} Toute fonction dérivable en $a$ est continue en $a$. \emph{La réciproque est fausse.}
\end{itemize}
\end{propriete}

\begin{demo}
Montrons que dérivable en $a$ implique continue en $a$. Pour $x \neq a$ au voisinage de $a$ :
\[
  f(x) - f(a) = \frac{f(x) - f(a)}{x - a} \times (x - a).
\]
Quand $x \to a$, le premier facteur tend vers $f'(a)$ (fini) et le second vers $0$ : par produit des limites, $f(x) - f(a) \to 0$, c'est-à-dire $\lim\limits_{x \to a} f(x) = f(a)$, ce qui est la continuité en $a$.
\end{demo}

\begin{exemple}[Calcul direct du nombre dérivé]
Montrons, à partir de la définition, que $f(x) = x^2$ est dérivable en $3$ et calculons $f'(3)$.
\[
  \frac{f(3 + h) - f(3)}{h} = \frac{(3+h)^2 - 9}{h} = \frac{9 + 6h + h^2 - 9}{h} = \frac{6h + h^2}{h} = 6 + h
  \xrightarrow[h \to 0]{} 6.
\]
Donc $f$ est dérivable en $3$ et $f'(3) = 6$.
\end{exemple}

\begin{exemple}[Point anguleux]
Étudions la dérivabilité de $f(x) = |x|$ en $0$.
\[
  \lim_{x \to 0^+} \frac{|x| - 0}{x - 0} = \lim_{x \to 0^+} \frac{x}{x} = 1,
  \qquad
  \lim_{x \to 0^-} \frac{|x|}{x} = \lim_{x \to 0^-} \frac{-x}{x} = -1.
\]
$f'_d(0) = 1 \neq f'_g(0) = -1$ : $f$ n'est pas dérivable en $0$, bien qu'elle y soit continue. La courbe présente un point anguleux à l'origine, avec deux demi-tangentes de pentes $1$ et $-1$.
\end{exemple}

\begin{exemple}[Tangente verticale]
Étudions la dérivabilité de $f(x) = \sqrt{x}$ en $0$.
\[
  \frac{f(x) - f(0)}{x - 0} = \frac{\sqrt{x}}{x} = \frac{1}{\sqrt{x}} \xrightarrow[x \to 0^+]{} +\infty.
\]
$f$ n'est pas dérivable en $0$ (limite infinie) ; la courbe admet en $0$ une \textbf{demi-tangente verticale} d'équation $x = 0$.
\end{exemple}

\begin{remarque}[Approximation affine locale]
Au voisinage de $a$, si $f$ est dérivable en $a$ : $f(x) \approx f(a) + f'(a)(x - a)$ pour $x$ proche de $a$. Utile pour estimer une valeur numérique sans calculatrice : par exemple, avec $f(x) = \sqrt{x}$, $f'(x) = \frac{1}{2\sqrt{x}}$, donc $f'(4) = \frac{1}{4}$, d'où
\[
  \sqrt{4{,}02} \approx f(4) + f'(4)\times 0{,}02 = 2 + \frac{0{,}02}{4} = 2{,}005.
\]
\end{remarque}

% ------------------- SECTION 2 -------------------
\section{Calcul des dérivées : toutes les formules}

\begin{propriete}[Dérivées des fonctions usuelles]
\begin{center}
\renewcommand{\arraystretch}{1.5}
\sffamily\small
\begin{tabular}{>{\centering\arraybackslash}p{3.6cm} >{\centering\arraybackslash}p{4.4cm} >{\centering\arraybackslash}p{6.4cm}}
\rowcolor{qtealdark} \color{white}\bfseries Fonction & \color{white}\bfseries Dérivée & \color{white}\bfseries Domaine de dérivabilité \\
\rowcolor{tealpale} $x^n$ ($n \in \mathbb{Z}$) & $n\,x^{n-1}$ & $\mathbb{R}$ (ou $\mathbb{R}^*$ si $n < 0$) \\
$\sqrt{x}$ & $\dfrac{1}{2\sqrt{x}}$ & $]0\,;+\infty[$ \\
\rowcolor{tealpale} $\sqrt[n]{x}$ & $\dfrac{1}{n}x^{\frac{1}{n}-1}$ & $]0\,;+\infty[$ \\
$\sin x$ & $\cos x$ & $\mathbb{R}$ \\
\rowcolor{tealpale} $\cos x$ & $-\sin x$ & $\mathbb{R}$ \\
$\tan x$ & $1 + \tan^2 x = \dfrac{1}{\cos^2 x}$ & $x \neq \frac{\pi}{2} + k\pi$ \\
\end{tabular}
\end{center}
\end{propriete}

\begin{propriete}[Opérations sur les dérivées]
Si $u$ et $v$ sont dérivables sur $I$ (avec $\lambda \in \mathbb{R}$) :
\[
  (u + v)' = u' + v', \qquad
  (\lambda u)' = \lambda u', \qquad
  (uv)' = u'v + uv',
\]
\[
  \left(\frac{u}{v}\right)' = \frac{u'v - uv'}{v^2} \ (v \neq 0), \qquad
  \left(\frac{1}{v}\right)' = \frac{-v'}{v^2}.
\]
\end{propriete}

\begin{propriete}[Dérivée d'une fonction composée : toutes les formes]
Si $g$ est dérivable en $f(a)$ et $f$ dérivable en $a$, alors $g \circ f$ est dérivable en $a$ et
\[
  (g \circ f)'(a) = f'(a) \times g'\bigl(f(a)\bigr).
\]
Formes composées usuelles, pour $u$ dérivable sur $I$ :
\[
  (u^n)' = n\,u'\,u^{n-1} \ (n \in \mathbb{Z}), \qquad
  \left(\sqrt{u}\right)' = \frac{u'}{2\sqrt{u}} \ (u > 0), \qquad
  \left(\sqrt[n]{u}\right)' = \frac{u'}{n\,u^{\frac{n-1}{n}}} \ (u > 0),
\]
\[
  \bigl(\sin u\bigr)' = u'\cos u, \qquad
  \bigl(\cos u\bigr)' = -u'\sin u, \qquad
  \bigl(\tan u\bigr)' = u'\bigl(1 + \tan^2 u\bigr) = \frac{u'}{\cos^2 u}.
\]
\end{propriete}

\begin{demo}
Démontrons $(u^2)' = 2u'u$ à partir de la règle du produit : en écrivant $u^2 = u \times u$,
\[
  (u \times u)' = u' \times u + u \times u' = 2u\,u'.
\]
Le cas général $(u^n)'$ se démontre par récurrence sur $n$ en utilisant cette même règle du produit à chaque étape.
\end{demo}

\begin{propriete}[Dérivée de la fonction réciproque]
Soit $f$ continue, strictement monotone sur $I$, dérivable en $a \in I$ avec $f'(a) \neq 0$. Alors $f^{-1}$ est dérivable en $b = f(a)$ et
\[
  \left(f^{-1}\right)'(b) = \frac{1}{f'\bigl(f^{-1}(b)\bigr)}.
\]
\textbf{Cas où $f'(a) = 0$} : $f^{-1}$ n'est \textbf{pas} dérivable en $b = f(a)$ (la courbe de $f^{-1}$ admet une tangente verticale en ce point).

\smallskip
Application aux puissances rationnelles : pour $x > 0$ et $r \in \mathbb{Q}^*$,
\[
  \left(x^{r}\right)' = r\,x^{r-1}.
\]
\end{propriete}

\begin{demo}
Partons de $f\bigl(f^{-1}(y)\bigr) = y$ pour tout $y \in J$, et dérivons formellement les deux membres (composée), en admettant que $f^{-1}$ est dérivable en $b$ :
\[
  \bigl(f^{-1}\bigr)'(b) \times f'\bigl(f^{-1}(b)\bigr) = 1
  \quad\Longrightarrow\quad
  \bigl(f^{-1}\bigr)'(b) = \frac{1}{f'\bigl(f^{-1}(b)\bigr)}.
\]
\end{demo}

\begin{exemple}[Dérivée d'une réciproque -- calcul complet]
Soit $f(x) = x^3 + x$ sur $\mathbb{R}$, strictement croissante ($f'(x) = 3x^2 + 1 > 0$), donc bijective de $\mathbb{R}$ sur $\mathbb{R}$. Calculons $(f^{-1})'(2)$ sans expliciter $f^{-1}$.

\smallskip
On cherche $a$ tel que $f(a) = 2$ : on constate $f(1) = 1 + 1 = 2$, donc $f^{-1}(2) = 1$. Comme $f'(1) = 3 + 1 = 4 \neq 0$ :
\[
  (f^{-1})'(2) = \frac{1}{f'(1)} = \frac{1}{4}.
\]
\end{exemple}

\begin{exemple}[Dérivées composées combinées]
Dérivons $f(x) = \sqrt{3x^2 + 1}$ et $g(x) = \sin(2x^3 - x)$ sur $\mathbb{R}$.
\[
  f'(x) = \frac{(3x^2+1)'}{2\sqrt{3x^2+1}} = \frac{6x}{2\sqrt{3x^2+1}} = \frac{3x}{\sqrt{3x^2+1}},
\]
\[
  g'(x) = (6x^2 - 1)\cos(2x^3 - x).
\]
\end{exemple}

\begin{exemple}[Tangente à une courbe]
Écrivons l'équation de la tangente à $\mathcal{C}_f$ en $x_0 = 1$, où $f(x) = x^3 - 2x + 1$.
\[
  f(1) = 1 - 2 + 1 = 0, \qquad f'(x) = 3x^2 - 2, \qquad f'(1) = 1.
\]
Tangente : $y = f'(1)(x - 1) + f(1) = 1\times(x - 1) + 0 = x - 1$.
\end{exemple}

% ------------------- SECTION 3 -------------------
\section{Théorème de Rolle et accroissements finis}

\begin{propriete}[Théorème de Rolle]
Soit $f$ \textbf{continue} sur $[a\,;b]$, \textbf{dérivable} sur $]a\,;b[$, avec $f(a) = f(b)$. Alors il existe (au moins) un $c \in\, ]a\,;b[$ tel que
\[
  f'(c) = 0.
\]
\end{propriete}

\begin{demo}
$f$ est continue sur le segment $[a\,;b]$ : elle est donc bornée et atteint son minimum $m$ et son maximum $M$ (propriété admise, vue au chapitre précédent).

\smallskip
\textbf{Cas 1} : $m = M$. Alors $f$ est constante sur $[a\,;b]$, donc $f'(c) = 0$ pour tout $c \in\, ]a\,;b[$.

\smallskip
\textbf{Cas 2} : $m \neq M$. Comme $f(a) = f(b)$, l'un au moins des deux extrema $m$ ou $M$ est atteint en un point $c \in\, ]a\,;b[$ (pas aux bornes). En ce point $c$, $f$ admet un \textbf{extremum local}, et $f$ est dérivable en $c$ : par la condition nécessaire d'extremum (voir section suivante), $f'(c) = 0$.
\end{demo}

\begin{propriete}[Théorème des accroissements finis (TAF)]
Soit $f$ \textbf{continue} sur $[a\,;b]$ et \textbf{dérivable} sur $]a\,;b[$. Alors il existe $c \in\, ]a\,;b[$ tel que
\[
  f(b) - f(a) = f'(c)\,(b - a).
\]
\textbf{Inégalité des accroissements finis} : si de plus $m \leq f'(x) \leq M$ sur $]a\,;b[$, alors
\[
  m\,(b - a) \;\leq\; f(b) - f(a) \;\leq\; M\,(b - a),
\]
et si $|f'(x)| \leq k$ sur l'intervalle, alors $|f(b) - f(a)| \leq k\,|b - a|$.
\end{propriete}

\begin{demo}
Le TAF se déduit de Rolle en introduisant la fonction auxiliaire
\[
  \varphi(x) = f(x) - f(a) - \frac{f(b) - f(a)}{b - a}(x - a).
\]
$\varphi$ est continue sur $[a\,;b]$, dérivable sur $]a\,;b[$, et $\varphi(a) = 0$, $\varphi(b) = f(b) - f(a) - \bigl(f(b) - f(a)\bigr) = 0$. D'après Rolle, il existe $c \in\, ]a\,;b[$ tel que $\varphi'(c) = 0$, c'est-à-dire
\[
  f'(c) - \frac{f(b) - f(a)}{b - a} = 0 \quad\Longrightarrow\quad f(b) - f(a) = f'(c)(b - a).
\]
L'inégalité s'obtient en encadrant $f'(c)$ par $m$ et $M$, puis en multipliant par $(b - a) > 0$.
\end{demo}

\begin{exemple}[Encadrement classique via TAF]
Montrons que pour tout $x > 0$ : $\ \sin x \leq x$.

\smallskip
Posons $f(t) = \sin t$ sur $[0\,;x]$ : $f$ est continue sur $[0\,;x]$, dérivable sur $]0\,;x[$ et $f'(t) = \cos t \leq 1$. L'inégalité des accroissements finis avec $M = 1$ donne
\[
  \sin x - \sin 0 \leq 1 \times (x - 0), \qquad\text{d'où}\qquad \sin x \leq x.
\]
\end{exemple}

\begin{exemple}[Encadrement d'une racine]
Montrons que pour tout $x > 1$ : $\ \dfrac{1}{2\sqrt{x}} < \sqrt{x} - \sqrt{x - 1} < \dfrac{1}{2\sqrt{x-1}}$.

\smallskip
Appliquons le TAF à $f(t) = \sqrt{t}$ sur $[x-1\,;x]$ : il existe $c \in\, ]x-1\,;x[$ tel que $\sqrt{x} - \sqrt{x-1} = f'(c) = \dfrac{1}{2\sqrt{c}}$. Comme $x - 1 < c < x$, la fonction $t \mapsto \frac{1}{2\sqrt{t}}$ étant décroissante :
\[
  \frac{1}{2\sqrt{x}} < \frac{1}{2\sqrt{c}} < \frac{1}{2\sqrt{x-1}},
\]
ce qui donne exactement l'encadrement annoncé.
\end{exemple}

\begin{remarque}[Utilisation type bac avec les suites récurrentes]
L'inégalité $|f(b) - f(a)| \leq k|b - a|$ sert dans les problèmes de suites $u_{n+1} = f(u_n)$ avec point fixe $\ell$ ($f(\ell) = \ell$) : si $|f'| \leq k < 1$ sur un intervalle stable contenant les $u_n$ et $\ell$, alors
\[
  |u_{n+1} - \ell| = |f(u_n) - f(\ell)| \leq k\,|u_n - \ell|,
\]
d'où, par récurrence immédiate, $|u_n - \ell| \leq k^n\,|u_0 - \ell| \to 0$ (car $0 \leq k < 1$) : ceci prouve la convergence de $(u_n)$ vers $\ell$, \textbf{sans} passer par monotonie + borne.
\end{remarque}

% ------------------- SECTION 4 -------------------
\section{Monotonie, extremums, concavité}

\begin{propriete}[Condition nécessaire d'extremum local]
Si $f$ est dérivable en $a$, intérieur à son domaine, et si $f$ admet un extremum local en $a$, alors $f'(a) = 0$.
\end{propriete}

\begin{remarque}[La réciproque est fausse]
$f'(a) = 0$ ne suffit \emph{pas} à garantir un extremum : pour $f(x) = x^3$, $f'(0) = 0$ mais $f$ est strictement croissante sur $\mathbb{R}$ (pas d'extremum en $0$, c'est un \textbf{point d'inflexion à tangente horizontale}). Il faut vérifier que $f'$ \textbf{change de signe} en $a$.
\end{remarque}

\begin{propriete}[Signe de la dérivée et variations]
Soit $f$ dérivable sur un intervalle $I$ :
\begin{itemize}[leftmargin=6mm, itemsep=0.5mm]
  \item $f' \geq 0$ sur $I$ (nulle en des points isolés) $\iff$ $f$ croissante sur $I$ ; $f' > 0$ sur $I$ $\Rightarrow$ $f$ strictement croissante ;
  \item $f' \leq 0$ sur $I$ $\iff$ $f$ décroissante sur $I$ ; $f' < 0$ sur $I$ $\Rightarrow$ $f$ strictement décroissante ;
  \item $f' = 0$ sur $I$ $\iff$ $f$ constante sur $I$.
\end{itemize}
\textbf{Extremum local} (condition suffisante) : si $f'$ s'annule en $a$ \textbf{en changeant de signe}, alors $f(a)$ est un extremum local : maximum si $f'$ passe de $+$ à $-$, minimum si $f'$ passe de $-$ à $+$.
\end{propriete}

\begin{demo}
($f' \geq 0 \Rightarrow f$ croissante) Soient $x_1 < x_2$ dans $I$. Par le TAF appliqué sur $[x_1\,;x_2]$, il existe $c \in\, ]x_1\,;x_2[$ tel que $f(x_2) - f(x_1) = f'(c)(x_2 - x_1)$. Comme $f'(c) \geq 0$ et $x_2 - x_1 > 0$, on a $f(x_2) - f(x_1) \geq 0$, donc $f(x_1) \leq f(x_2)$ : $f$ est croissante.
\end{demo}

\begin{definition}[Concavité et point d'inflexion]
Soit $f$ deux fois dérivable sur $I$.
\begin{itemize}[leftmargin=6mm, itemsep=0.5mm]
  \item Si $f'' \geq 0$ sur $I$, la courbe est \textbf{convexe} sur $I$ (elle est située au-dessus de chacune de ses tangentes).
  \item Si $f'' \leq 0$ sur $I$, la courbe est \textbf{concave} sur $I$ (au-dessous de ses tangentes).
  \item Un \textbf{point d'inflexion} est un point où $f''$ s'annule \textbf{en changeant de signe} : la courbe y traverse sa tangente.
\end{itemize}
\end{definition}

\begin{propriete}[Convexité et position par rapport à la tangente]
Si $f$ est convexe sur $I$, pour tout $a \in I$ et tout $x \in I$ :
\[
  f(x) \geq f(a) + f'(a)(x - a).
\]
(Autrement dit, la courbe est toujours au-dessus de sa tangente en tout point.) L'inégalité est renversée si $f$ est concave.
\end{propriete}

\begin{remarque}[Piège classique]
$f''(a) = 0$ ne suffit \emph{pas} pour un point d'inflexion : il faut le \emph{changement de signe} de $f''$. Contre-exemple : $f(x) = x^4$ en $0$ : $f''(x) = 12x^2$, donc $f''(0) = 0$, mais $f''(x) \geq 0$ partout (pas de changement de signe) : la courbe reste convexe, il n'y a \textbf{pas} d'inflexion en $0$.
\end{remarque}

\begin{exemple}[Extremum et convexité]
Étudions $f(x) = x^3 - 3x$ sur $\mathbb{R}$ : variations, extrema, concavité, inflexion.
\[
  f'(x) = 3x^2 - 3 = 3(x-1)(x+1), \qquad f''(x) = 6x.
\]
$f'$ s'annule en $-1$ et $1$, change de signe ($+,-,+$) : \textbf{maximum local} $f(-1) = 2$, \textbf{minimum local} $f(1) = -2$.

\smallskip
$f''$ s'annule en $0$, change de signe ($f'' < 0$ sur $]-\infty\,;0[$, $f'' > 0$ sur $]0\,;+\infty[$) : \textbf{point d'inflexion} en $x = 0$ (avec $f(0) = 0$), la courbe passant de concave à convexe.
\end{exemple}

% ------------------- SECTION 5 -------------------
\section{Branches infinies : classification complète}

\begin{propriete}[Les quatre types de branches infinies]
\begin{enumerate}[leftmargin=7mm, itemsep=0.5mm]
  \item \textbf{Asymptote verticale} : $\lim\limits_{x \to a} f(x) = \pm\infty$ (ou une seule limite latérale infinie) $\ \Rightarrow$ droite $x = a$.
  \item \textbf{Asymptote horizontale} : $\lim\limits_{x \to \pm\infty} f(x) = b$ (fini) $\ \Rightarrow$ droite $y = b$.
  \item \textbf{Asymptote oblique} : $\lim\limits_{x \to \pm\infty} \bigl[f(x) - (ax + b)\bigr] = 0$ $\ \Rightarrow$ droite $y = ax + b$ ($a \neq 0$).
  \item \textbf{Branches paraboliques} : quand $f(x) \to \pm\infty$ sans asymptote, on étudie $\dfrac{f(x)}{x}$ :
    \begin{itemize}[leftmargin=5mm]
      \item $\dfrac{f(x)}{x} \to \pm\infty$ : branche parabolique de \textbf{direction $(Oy)$} ;
      \item $\dfrac{f(x)}{x} \to 0$ : branche parabolique de \textbf{direction $(Ox)$} ;
      \item $\dfrac{f(x)}{x} \to a \neq 0$ (fini) : on étudie ensuite $f(x) - ax$ :
        \begin{itemize}[leftmargin=5mm]
          \item limite finie $b$ $\Rightarrow$ \textbf{asymptote oblique} $y = ax + b$ (retour au cas 3) ;
          \item limite infinie $\Rightarrow$ \textbf{branche parabolique de direction} $y = ax$.
        \end{itemize}
    \end{itemize}
\end{enumerate}
\end{propriete}

\begin{propriete}[Position de la courbe par rapport à l'asymptote oblique]
Si $\mathcal{C}_f$ admet l'asymptote $y = ax + b$ en $+\infty$ (ou $-\infty$), on étudie le signe de $f(x) - (ax + b)$ :
\begin{itemize}[leftmargin=6mm, itemsep=0.5mm]
  \item $f(x) - (ax+b) > 0$ : la courbe est \textbf{au-dessus} de l'asymptote ;
  \item $f(x) - (ax+b) < 0$ : la courbe est \textbf{au-dessous} de l'asymptote ;
  \item si ce signe change en un point, la courbe \textbf{traverse} son asymptote en ce point.
\end{itemize}
\end{propriete}

\begin{exemple}[Étude guidée complète]
Étudions $f(x) = \dfrac{x^2 + 1}{x}$ sur $D_f = \mathbb{R}^*$.

\smallskip
\textbf{Écriture simplifiée} : $f(x) = x + \dfrac{1}{x}$. $f$ est impaire ($f(-x) = -f(x)$) : on étudie sur $]0\,;+\infty[$ puis on complète par symétrie par rapport à l'origine.

\smallskip
\textbf{Limites} : $\lim\limits_{x \to 0^+} f(x) = +\infty$ : \textbf{asymptote verticale} $x = 0$. En $+\infty$ : $f(x) \to +\infty$, et $f(x) - x = \frac{1}{x} \to 0$ : \textbf{asymptote oblique} $y = x$. Signe de $f(x) - x = \frac{1}{x} > 0$ sur $]0\,;+\infty[$ : la courbe est \textbf{au-dessus} de l'asymptote.

\smallskip
\textbf{Dérivée} : $f'(x) = 1 - \dfrac{1}{x^2} = \dfrac{x^2 - 1}{x^2}$, du signe de $x^2 - 1$. Sur $]0\,;+\infty[$ : $f' < 0$ sur $]0\,;1[$, $f' > 0$ sur $]1\,;+\infty[$ : \textbf{minimum local} en $x = 1$, $f(1) = 2$.

\smallskip
\textbf{Conclusion} : sur $]0\,;+\infty[$, $f$ décroît de $+\infty$ à $2$ puis croît vers $+\infty$ ; asymptotes $x = 0$ et $y = x$ ; par imparité, sur $]-\infty\,;0[$ : maximum local $f(-1) = -2$, mêmes asymptotes.
\end{exemple}

\subsection{Plan d'étude complet -- méthode à suivre}

\begin{remarque}[Les six étapes]
\begin{enumerate}[leftmargin=7mm, itemsep=0.5mm]
  \item \textbf{Domaine} $D_f$ ; parité, périodicité éventuelles (pour réduire l'intervalle d'étude).
  \item \textbf{Limites} aux bornes de $D_f$ ; classification des asymptotes et branches infinies.
  \item \textbf{Dérivabilité}, calcul de $f'$, étude de son signe.
  \item \textbf{Tableau de variations} complet (flèches, valeurs et limites en chaque borne).
  \item \textbf{Concavité}, points d'inflexion (si demandé), tangentes remarquables (points à tangente horizontale, verticale).
  \item \textbf{Tracé} de la courbe : asymptotes en pointillés, extremums, points d'inflexion, quelques valeurs numériques.
\end{enumerate}
\end{remarque}

% ------------------- SECTION 6 -------------------
\section{Exercices corrigés -- niveau examen national SM}

\begin{exemple}[Exercice 1 -- dérivabilité et paramètre]
Soit $f$ définie par $f(x) = \sqrt{x^2 + 1}$ si $x \leq 0$ et $f(x) = ax + 1$ si $x > 0$. Déterminer $a$ pour que $f$ soit dérivable en $0$.

\smallskip
\textbf{Corrigé.} \emph{Continuité d'abord} : $f(0) = 1$ (branche de gauche) ; $\lim\limits_{x \to 0^+} f(x) = 1$ (branche de droite, quel que soit $a$) : continuité automatique.

\smallskip
\emph{Dérivées à gauche et à droite} : pour $x \leq 0$, $f'(x) = \dfrac{x}{\sqrt{x^2+1}}$, donc $f'_g(0) = 0$. Pour $x > 0$, $f'(x) = a$, donc $f'_d(0) = a$.

\smallskip
$f$ dérivable en $0$ $\iff f'_g(0) = f'_d(0) \iff a = 0$.
\end{exemple}

\begin{exemple}[Exercice 2 -- tangente et position relative]
Soit $f(x) = x^3 - 3x^2 + 2$. Écrire l'équation de la tangente $T$ au point d'abscisse $1$, et étudier la position de $\mathcal{C}_f$ par rapport à $T$.

\smallskip
\textbf{Corrigé.} $f(1) = 1 - 3 + 2 = 0$ ; $f'(x) = 3x^2 - 6x$, $f'(1) = -3$. Tangente : $y = -3(x-1) + 0 = -3x + 3$.

\smallskip
Étudions $g(x) = f(x) - (-3x + 3) = x^3 - 3x^2 + 3x - 1$. On reconnaît l'identité remarquable $g(x) = (x - 1)^3$ (vérification : $(x-1)^3 = x^3 - 3x^2 + 3x - 1$ \checkmark).

\smallskip
Le signe de $(x-1)^3$ est celui de $(x - 1)$ : $g(x) < 0$ pour $x < 1$, $g(x) = 0$ en $x=1$, $g(x) > 0$ pour $x > 1$.

\smallskip
\textbf{Conclusion} : $\mathcal{C}_f$ est en dessous de $T$ pour $x < 1$, au-dessus pour $x > 1$ : la courbe \textbf{traverse} sa tangente en $x=1$, qui est donc un point d'inflexion (vérification : $f''(x) = 6x - 6$, s'annule bien en $1$ en changeant de signe \checkmark).
\end{exemple}

\begin{exemple}[Exercice 3 -- Rolle appliqué à une dérivée]
Soit $f(x) = (x-1)(x-2)(x-3)$. Montrer, sans calculer $f'$, que $f'$ admet au moins deux racines réelles distinctes.

\smallskip
\textbf{Corrigé.} $f$ est un polynôme, donc continue et dérivable sur $\mathbb{R}$. On a $f(1) = f(2) = f(3) = 0$.

\smallskip
Sur $[1\,;2]$ : $f$ continue, dérivable, $f(1) = f(2)$ : par Rolle, il existe $c_1 \in\, ]1\,;2[$ tel que $f'(c_1) = 0$.

\smallskip
Sur $[2\,;3]$ : de même, il existe $c_2 \in\, ]2\,;3[$ tel que $f'(c_2) = 0$.

\smallskip
Comme $]1\,;2[$ et $]2\,;3[$ sont disjoints, $c_1 \neq c_2$ : $f'$ admet (au moins) deux racines réelles distinctes.
\end{exemple}

\begin{exemple}[Exercice 4 -- suite récurrente via l'inégalité des accroissements finis]
Soit $f(x) = \sqrt{x + 6}$, et $(u_n)$ définie par $u_0 = 3$, $u_{n+1} = f(u_n)$. On admet que $u_n \in [2\,;3]$ pour tout $n$. Montrer que $|u_{n+1} - 3| \leq \frac{1}{4}|u_n - 3|$, puis en déduire $\lim u_n$.

\smallskip
\textbf{Corrigé.} $f(3) = \sqrt{9} = 3$ : $3$ est point fixe de $f$. Sur $[2\,;3]$, $f'(x) = \dfrac{1}{2\sqrt{x+6}}$ ; comme $x + 6 \in [8\,;9]$, $\sqrt{x+6} \geq \sqrt{8} = 2\sqrt{2}$, donc $f'(x) \leq \dfrac{1}{4\sqrt{2}} < \dfrac{1}{4}$.

\smallskip
Par l'inégalité des accroissements finis entre $u_n$ et $3$ (tous deux dans $[2\,;3]$, intervalle stable) :
\[
  |u_{n+1} - 3| = |f(u_n) - f(3)| \leq \frac{1}{4}|u_n - 3|.
\]
Par récurrence : $|u_n - 3| \leq \left(\frac{1}{4}\right)^n |u_0 - 3| = \left(\frac{1}{4}\right)^n \to 0$. Par encadrement, $u_n - 3 \to 0$, donc $\boxed{\lim u_n = 3}$.
\end{exemple}

\begin{exemple}[Exercice 5 -- convexité et inégalité]
Montrer que pour tout $x \geq 0$ : $\sqrt{1 + x} \leq 1 + \dfrac{x}{2}$.

\smallskip
\textbf{Corrigé.} Posons $f(t) = \sqrt{1+t}$ sur $[0\,;+\infty[$. $f'(t) = \dfrac{1}{2\sqrt{1+t}}$, $f''(t) = -\dfrac{1}{4(1+t)^{3/2}} < 0$ : $f$ est \textbf{concave} sur $[0\,;+\infty[$.

\smallskip
Une fonction concave est en dessous de chacune de ses tangentes. La tangente en $t = 0$ : $f(0) = 1$, $f'(0) = \frac{1}{2}$, donc $y = 1 + \frac{t}{2}$. Par concavité :
\[
  f(t) \leq 1 + \frac{t}{2} \qquad \text{pour tout } t \geq 0,
\]
ce qui est exactement l'inégalité demandée (avec $t = x$).
\end{exemple}

\begin{exemple}[Exercice 6 -- étude complète avec branches infinies]
Étudier $f(x) = \dfrac{x^2 - x + 1}{x - 1}$ sur $D_f = \mathbb{R} \setminus \{1\}$ : asymptotes, variations, tracé.

\smallskip
\textbf{Corrigé.} \emph{Écriture} : division euclidienne donne $f(x) = x + \dfrac{1}{x-1}$.

\smallskip
\emph{Limites} : $\lim\limits_{x \to 1^{\pm}} f(x) = \pm\infty$ : asymptote verticale $x = 1$. En $\pm\infty$ : $f(x) - x = \frac{1}{x-1} \to 0$ : asymptote oblique $y = x$. Signe de $\frac{1}{x-1}$ : positif si $x > 1$ (courbe au-dessus), négatif si $x < 1$ (courbe en dessous).

\smallskip
\emph{Dérivée} : $f'(x) = 1 - \dfrac{1}{(x-1)^2} = \dfrac{(x-1)^2 - 1}{(x-1)^2} = \dfrac{x(x-2)}{(x-1)^2}$.

\smallskip
Signe de $f'$ = signe de $x(x-2)$ (dénominateur toujours positif) : $f' > 0$ sur $]-\infty\,;0[$, $f' < 0$ sur $]0\,;1[\cup]1\,;2[$, $f' > 0$ sur $]2\,;+\infty[$.

\smallskip
\textbf{Conclusion} : maximum local $f(0) = -1$, minimum local $f(2) = 3$ ; asymptotes $x=1$ et $y=x$.
\end{exemple}

% ------------------- SECTION 7 -------------------
\section{Méthodes et pièges : la boîte à outils de l'examen}

\subsection{Ce qu'on te demande, ce que tu fais}

\begin{center}
\renewcommand{\arraystretch}{1.45}
\sffamily\small
\begin{tabular}{>{\raggedright\arraybackslash}p{7.2cm} >{\raggedright\arraybackslash}p{8cm}}
\rowcolor{qtealdark} \color{white}\bfseries Ce qu'on te demande & \color{white}\bfseries Ton réflexe \\
\rowcolor{tealpale} Étudier la dérivabilité en $a$ & Calculer $f'_g(a)$ et $f'_d(a)$ séparément, comparer. \\
Écrire l'équation d'une tangente & $y = f'(x_0)(x - x_0) + f(x_0)$. \\
\rowcolor{tealpale} Position courbe / tangente ou asymptote & Étudier le signe de $f(x) - (\text{droite})$. \\
Montrer qu'une dérivée s'annule $k$ fois & Rolle sur chacun des $k$ sous-intervalles où $f$ prend la même valeur. \\
\rowcolor{tealpale} Prouver la convergence d'une suite $u_{n+1}=f(u_n)$ & Inégalité des accroissements finis : $|f'| \leq k < 1$ sur un intervalle stable. \\
Trouver les extrema locaux & Signe de $f'$ ; vérifier le \textbf{changement de signe}, pas juste $f'=0$. \\
\rowcolor{tealpale} Trouver un point d'inflexion & Signe de $f''$ ; vérifier le \textbf{changement de signe}. \\
Prouver une inégalité $f(x) \leq $ droite & Étudier la concavité de $f$ (ou TAF). \\
\rowcolor{tealpale} Classifier une branche infinie & Étudier $\frac{f(x)}{x}$ puis $f(x) - ax$ selon l'arbre de décision. \\
Dériver une composée & Identifier $u$, appliquer la formule adaptée ($u^n$, $\sqrt{u}$, $\sin u$, etc.). \\
\end{tabular}
\end{center}

\subsection{Les pièges qui coûtent des points}

\begin{remarque}[Erreurs relevées chaque année par les correcteurs]
\begin{itemize}[leftmargin=6mm, itemsep=0.5mm]
  \item Confondre $f'(a) = 0$ avec « extremum en $a$ » : il faut le \textbf{changement de signe} (contre-exemple $x^3$ en $0$).
  \item Confondre $f''(a) = 0$ avec « point d'inflexion en $a$ » : même exigence de changement de signe (contre-exemple $x^4$ en $0$).
  \item Oublier de vérifier la \textbf{continuité} avant d'étudier la dérivabilité en un point de raccordement.
  \item Appliquer Rolle sans vérifier $f(a) = f(b)$, ou le TAF en oubliant l'hypothèse de continuité sur le \textbf{segment fermé}.
  \item Écrire $(f^{-1})'(b) = \frac{1}{f'(b)}$ au lieu de $\frac{1}{f'(f^{-1}(b))}$.
  \item Utiliser l'inégalité des accroissements finis sur un intervalle où $f'$ n'est \textbf{pas} bornée par $k < 1$.
  \item Oublier le facteur $u'$ en dérivant une composée : $(\sin(x^2))' = 2x\cos(x^2)$, \textbf{pas} $\cos(x^2)$.
  \item Ne pas distinguer asymptote oblique et branche parabolique : toujours vérifier que $f(x) - ax$ tend vers une limite \textbf{finie}.
  \item Conclure à la dérivabilité d'une fonction définie par morceaux sans comparer $f'_g$ et $f'_d$ séparément.
\end{itemize}
\end{remarque}

% ------------------- RESUME -------------------
\vspace{4mm}
\begin{tcolorbox}[
  enhanced, boxrule=0pt, arc=2mm,
  interior style={left color=qtealdark, right color=qteal},
  left=6mm, right=6mm, top=4mm, bottom=4mm,
  fontupper=\color{white}\sffamily,
  title={\sffamily\bfseries\color{white}\ding{72}~~À retenir},
  colbacktitle=qtealdark, coltitle=white, boxrule=0pt
]
\begin{itemize}[leftmargin=6mm, label={\color{qamber}\ding{212}}, itemsep=1mm]
  \item $f'(a) = \lim\limits_{h \to 0}\frac{f(a+h)-f(a)}{h}$ ; tangente $y = f'(a)(x-a)+f(a)$ ; dérivable $\Rightarrow$ continue (jamais l'inverse).
  \item Formules composées : $(u^n)' = nu'u^{n-1}$, $(\sqrt{u})' = \frac{u'}{2\sqrt{u}}$, $(\sin u)'=u'\cos u$ ; réciproque : $(f^{-1})'(b) = \frac{1}{f'(f^{-1}(b))}$.
  \item Rolle : $f(a)=f(b) \Rightarrow \exists c, f'(c)=0$ ; TAF : $f(b)-f(a)=f'(c)(b-a)$ ; $|f'|\leq k \Rightarrow |f(b)-f(a)|\leq k|b-a|$ (clé pour les suites).
  \item Signe de $f'$ $\to$ monotonie (démontré via TAF) ; extremum = changement de signe de $f'$ ; inflexion = changement de signe de $f''$.
  \item Branches infinies : verticale/horizontale/oblique/parabolique, via l'étude de $f(x)/x$ puis $f(x)-ax$.
  \item Convexe $\Rightarrow$ courbe au-dessus des tangentes ; outil puissant pour démontrer des inégalités.
\end{itemize}
\end{tcolorbox}

\end{document}
